% ════════════════════════════════════════════════════════════
%  CHAPITRE 3 — IMPLÉMENTATION ET TESTS
%  Module : Réclamations & Justifications
%  Université Ibn Khaldoun — Tiaret
%  Groupe 5 — ISIL
%  Membres : Harbane Batoul · Hafsa Aicha Hayat · Hamani Ikram
% ════════════════════════════════════════════════════════════

% ── Préambule (à inclure dans le document principal) ────────
% \documentclass[12pt,a4paper]{report}
% \usepackage[utf8]{inputenc}
% \usepackage[T1]{fontenc}
% \usepackage[french]{babel}
% \usepackage{lmodern}
% \usepackage[top=2.5cm,bottom=2.5cm,left=3cm,right=2.5cm]{geometry}
% \usepackage{xcolor}
% \usepackage{tikz}
% \usetikzlibrary{shapes.geometric,arrows.meta,positioning}
% \usepackage{pgfplots}
% \pgfplotsset{compat=1.18}
% \usepackage{listings}
% \usepackage{booktabs}
% \usepackage{tabularx}
% \usepackage{float}
% \usepackage{caption}
% \usepackage[most]{tcolorbox}
% \usepackage{setspace}
% \onehalfspacing
%
% \definecolor{IBNBlue}{RGB}{30,80,160}
% \definecolor{IBNGold}{RGB}{180,140,30}
% \definecolor{LightBlue}{RGB}{235,242,255}
% \definecolor{AccentGray}{RGB}{100,100,110}
% ────────────────────────────────────────────────────────────

\chapter{Implémentation et Tests}

\section{Introduction}

Ce chapitre présente l'implémentation concrète du module
\textbf{Réclamations \& Justifications} du Système de Gestion Universitaire
Intégré d'Ibn Khaldoun. Après avoir défini l'architecture et les besoins dans
les chapitres précédents, nous décrivons ici les choix techniques effectués,
le code développé côté backend et frontend, puis les résultats des tests de
validation fonctionnelle.

% ────────────────────────────────────────────────────────────
\section{Environnement de Développement}
% ────────────────────────────────────────────────────────────

\begin{table}[H]
\centering
\caption{Environnement de développement — Module Réclamations \& Justifications}
\begin{tabular}{@{}ll@{}}
\toprule
\textbf{Composant}       & \textbf{Technologie / Version} \\
\midrule
Langage backend          & TypeScript 5.x                 \\
Runtime                  & Node.js 20 LTS                 \\
Framework API            & Express.js 4.18                \\
ORM                      & Prisma 5.x                     \\
Base de données          & PostgreSQL 16                  \\
Framework frontend       & React.js 18                    \\
Styles                   & Tailwind CSS 3.x               \\
Client HTTP              & Axios                          \\
Authentification         & JSON Web Token (JWT)           \\
Notifications            & Système interne (Prisma \texttt{Notification}) \\
\bottomrule
\end{tabular}
\end{table}

% ────────────────────────────────────────────────────────────
\section{Architecture du Module}
% ────────────────────────────────────────────────────────────

Le module suit l'architecture globale du SGUI~: séparation stricte entre
\textbf{backend RESTful} (Node.js/Express) et \textbf{frontend SPA}
(React.js), communicant via des requêtes HTTP/JSON authentifiées par JWT.

\begin{figure}[H]
\centering
\begin{tikzpicture}[
  box/.style={draw=IBNBlue,thick,rounded corners=5pt,
              minimum width=3.0cm,minimum height=0.9cm,
              fill=LightBlue,font=\small\bfseries,align=center},
  arrow/.style={->,>=Stealth,thick,IBNBlue}
]
\node[box] (FE)   at (0,   0) {Frontend\\React.js};
\node[box] (MW)   at (4,   0) {Middleware\\requireAuth};
\node[box] (CTRL) at (8,   0) {Controller\\requests};
\node[box] (SVC)  at (8,  -2) {Service\\notifications};
\node[box] (DB)   at (4,  -2) {PostgreSQL\\Prisma ORM};

\draw[arrow] (FE)   -- node[above,font=\scriptsize]{HTTP + JWT} (MW);
\draw[arrow] (MW)   -- node[above,font=\scriptsize]{req.user}   (CTRL);
\draw[arrow] (CTRL) -- node[right,font=\scriptsize]{notify}     (SVC);
\draw[arrow] (CTRL) -- node[above,font=\scriptsize]{Prisma}     (DB);
\draw[arrow] (SVC)  -- (DB);
\end{tikzpicture}
\caption{Architecture interne du module Réclamations \& Justifications}
\label{fig:rj_arch}
\end{figure}

% ────────────────────────────────────────────────────────────
\section{Modèle de Données}
% ────────────────────────────────────────────────────────────

\subsection{Schéma Prisma}

Le modèle de données implémente les deux entités identifiées lors de la
conception~: \texttt{Reclamation} pour les étudiants connectés et
\texttt{GuestReclamation} pour les visiteurs sans compte.

\begin{lstlisting}[language=SQL,
  caption=schema.prisma — Modèles Reclamation et GuestReclamation]
model Reclamation {
  id                 Int       @id @default(autoincrement())
  type               String    // "reclamation" | "justification"
  sujet              String
  message            String
  status             String    @default("submitted")
  reponse_admin      String?
  reponse_admin_par  String?
  reponse_admin_date DateTime?
  dateAbsence        DateTime?
  createdAt          DateTime  @default(now())
  updatedAt          DateTime  @updatedAt

  userId       String
  user         User   @relation(fields:[userId], references:[id])
  traiteParId  String?
  traiteParUser User? @relation("TraitePar",
                        fields:[traiteParId], references:[id])
  notifications Notification[]
}

model GuestReclamation {
  id        Int      @id @default(autoincrement())
  prenom    String
  nom       String
  email     String
  sujet     String
  message   String
  status    String   @default("submitted")
  createdAt DateTime @default(now())
}
\end{lstlisting}

\subsection{Cycle de Vie d'une Demande}

\begin{figure}[H]
\centering
\begin{tikzpicture}[
  state/.style={draw=IBNBlue,thick,rounded corners=18pt,
                minimum width=2.8cm,minimum height=0.8cm,
                font=\small\bfseries},
  arrow/.style={->,>=Stealth,thick,IBNBlue},
]
\node[state,fill=yellow!25]  (S) at (0,   0) {submitted};
\node[state,fill=blue!15]    (U) at (4,   0) {under\_review};
\node[state,fill=green!20]   (R) at (8,   0.8) {resolved};
\node[state,fill=red!15]     (J) at (8,  -0.8) {rejected};

\draw[arrow] (S) -- node[above,font=\scriptsize]{Admin ouvre} (U);
\draw[arrow] (U) -- node[above,font=\scriptsize]{Approuver}   (R);
\draw[arrow] (U) -- node[below,font=\scriptsize]{Rejeter}     (J);
\end{tikzpicture}
\caption{Machine d'états d'une réclamation / justification}
\label{fig:states}
\end{figure}

\begin{table}[H]
\centering
\caption{Correspondance statuts internes / libellés interface}
\begin{tabular}{@{}lll@{}}
\toprule
\textbf{Statut base de données} & \textbf{Libellé affiché} & \textbf{Couleur} \\
\midrule
\texttt{submitted}     & En attente & Jaune \\
\texttt{under\_review} & En cours   & Bleu  \\
\texttt{resolved}      & Acceptée   & Vert  \\
\texttt{rejected}      & Rejetée    & Rouge \\
\bottomrule
\end{tabular}
\end{table}

% ────────────────────────────────────────────────────────────
\section{Implémentation Backend}
% ────────────────────────────────────────────────────────────

\subsection{Organisation des Fichiers}

\begin{lstlisting}[language=bash,
  caption=Structure du module requests côté backend]
backend/src/modules/requests/
├── request.controller.ts   # Logique metier principale
├── request.routes.ts       # Definition des routes API
└── reclamation.service.ts  # Service de notification etudiant
\end{lstlisting}

\subsection{Routes API}

\begin{table}[H]
\centering
\caption{Routes API — Module Réclamations \& Justifications}
\begin{tabularx}{\linewidth}{@{}llLl@{}}
\toprule
\textbf{Méthode} & \textbf{Route}                    & \textbf{Action}                  & \textbf{Accès}  \\
\midrule
POST & \texttt{/reclamations}             & Soumettre une réclamation         & Étudiant        \\
POST & \texttt{/justifications}           & Soumettre une justification       & Étudiant        \\
GET  & \texttt{/my-reclamations}          & Voir ses réclamations             & Étudiant        \\
GET  & \texttt{/my-justifications}        & Voir ses justifications           & Étudiant        \\
GET  & \texttt{/admin/inbox}              & Boîte de réception globale        & Admin           \\
POST & \texttt{/admin/:id/decide}         & Approuver / Rejeter               & Admin           \\
POST & \texttt{/guest/reclamations}       & Soumettre (visiteur)              & Public          \\
\bottomrule
\end{tabularx}
\end{table}

\subsection{Soumission d'une Réclamation (Étudiant)}

\begin{lstlisting}[language=JavaScript,
  caption=request.controller.ts — Soumission d'une réclamation]
export const submitReclamation = async (req, res) => {
  const userId = req.user.sub;
  const { sujet, message, type } = req.body;

  if (!sujet?.trim() || !message?.trim())
    return res.status(400).json({
      error: { code: 'MISSING_FIELDS',
               message: 'Le sujet et le message sont obligatoires.' }
    });

  const reclamation = await prisma.reclamation.create({
    data: { userId, sujet, message,
            type: type || 'reclamation',
            status: 'submitted' }
  });

  res.status(201).json(reclamation);
};
\end{lstlisting}

\subsection{Soumission d'une Justification (Étudiant)}

\begin{lstlisting}[language=JavaScript,
  caption=request.controller.ts — Soumission d'une justification]
export const submitJustification = async (req, res) => {
  const userId = req.user.sub;
  const { sujet, message, dateAbsence } = req.body;

  if (!sujet?.trim() || !message?.trim() || !dateAbsence)
    return res.status(400).json({
      error: { code: 'MISSING_FIELDS',
               message: "Sujet, message et date d'absence sont obligatoires." }
    });

  const justification = await prisma.reclamation.create({
    data: { userId, sujet, message,
            type: 'justification',
            dateAbsence: new Date(dateAbsence),
            status: 'submitted' }
  });

  res.status(201).json(justification);
};
\end{lstlisting}

\subsection{Décision Administrative}

La fonction ci-dessous traite l'approbation ou le rejet d'une demande par
l'administrateur. Elle met à jour le statut, enregistre la réponse et déclenche
une notification automatique vers l'étudiant concerné.

\begin{lstlisting}[language=JavaScript,
  caption=request.controller.ts — Décision administrative]
// Supprime les prefixes internes [INFO_REQUEST] etc.
function stripDecisionPrefix(text) {
  return text
    .replace(/^\[(INFO_REQUEST|TEACHER_REVIEW|COUNCIL_REVIEW)\]\s*/i, '')
    .trim();
}

export const decideReclamation = async (req, res) => {
  const adminId        = req.user.sub;
  const { id }         = req.params;
  const { action, reponse } = req.body; // 'approve' | 'reject'

  if (!['approve', 'reject'].includes(action))
    return res.status(400).json({
      error: { code: 'INVALID_ACTION',
               message: 'Action invalide. Utiliser approve ou reject.' }
    });

  const existing = await prisma.reclamation.findUnique({
    where: { id: Number(id) }
  });
  if (!existing)
    return res.status(404).json({ error: { code: 'NOT_FOUND' } });

  if (['resolved', 'rejected'].includes(existing.status))
    return res.status(409).json({
      error: { code: 'ALREADY_REVIEWED',
               message: 'Cette demande a deja ete traitee.' }
    });

  const newStatus  = action === 'approve' ? 'resolved' : 'rejected';
  const cleanReply = reponse ? stripDecisionPrefix(reponse) : null;
  const now        = new Date();

  const updated = await prisma.reclamation.update({
    where: { id: Number(id) },
    data: {
      status:             newStatus,
      traiteParId:        adminId,
      reponse_admin:      cleanReply,
      reponse_admin_par:  adminId,
      reponse_admin_date: now,
      updatedAt:          now,
    }
  });

  await notifyStudentAboutReclamationDecision(
    existing.userId, updated, action
  );

  res.json({
    ...updated,
    adminResponse: { message: cleanReply, respondedBy: adminId, date: now }
  });
};
\end{lstlisting}

\subsection{Soumission par un Visiteur (Route Publique)}

\begin{lstlisting}[language=JavaScript,
  caption=request.controller.ts — Réclamation visiteur (sans authentification)]
export const submitGuestReclamation = async (req, res) => {
  const { prenom, nom, email, sujet, message } = req.body;

  if (!prenom?.trim() || !nom?.trim() || !email?.trim()
      || !sujet?.trim() || !message?.trim())
    return res.status(400).json({
      error: { code: 'MISSING_FIELDS',
               message: 'Tous les champs sont obligatoires.' }
    });

  const guest = await prisma.guestReclamation.create({
    data: { prenom, nom, email, sujet, message, status: 'submitted' }
  });

  res.status(201).json({
    message: 'Votre reclamation a bien ete recue.',
    id: guest.id
  });
};
\end{lstlisting}

\subsection{Service de Notification}

À chaque décision administrative, le système crée automatiquement une
notification dans la table \texttt{Notification} de l'étudiant concerné.

\begin{lstlisting}[language=JavaScript,
  caption=reclamation.service.ts — Notification automatique à l'étudiant]
export async function notifyStudentAboutReclamationDecision(
  studentId, reclamation, action
) {
  const label = action === 'approve' ? 'acceptee' : 'rejetee';

  await prisma.notification.create({
    data: {
      userId:  studentId,
      title:   `Reclamation ${label}`,
      message: `Votre reclamation "${reclamation.sujet}" `
             + `a ete ${label} par l'administration.`,
      type:    'reclamation_decision',
      read:    false,
    }
  });
}
\end{lstlisting}

% ────────────────────────────────────────────────────────────
\section{Implémentation Frontend}
% ────────────────────────────────────────────────────────────

\subsection{Organisation des Fichiers}

\begin{lstlisting}[language=bash,
  caption=Structure frontend — Module Réclamations \& Justifications]
frontend/src/
├── pages/
│   ├── RequestsPage.jsx        # Dashboard etudiant (liste + soumission)
│   ├── RequestDetailPage.jsx   # Detail d'une demande (conversation)
│   └── AdminRequestsPage.jsx   # Boite de reception admin
└── components/requests/
    └── RequestConversation.jsx # Fil de conversation admin <-> etudiant
\end{lstlisting}

\subsection{Normalisation des Données (Page Étudiant)}

\begin{lstlisting}[language=JavaScript,
  caption=RequestsPage.jsx — Normalisation du statut et de la réponse admin]
const STATUS_CONFIG = {
  pending:  { label: 'En attente', color: 'yellow' },
  reviewed: { label: 'En cours',   color: 'blue'   },
  accepted: { label: 'Acceptee',   color: 'green'  },
  rejected: { label: 'Rejetee',    color: 'red'    },
};

const DECISION_PREFIX_RE =
  /^\[(INFO_REQUEST|TEACHER_REVIEW|COUNCIL_REVIEW)\]\s*/i;

function cleanReplyText(text) {
  return text ? text.replace(DECISION_PREFIX_RE, '').trim() : '';
}

function normalizeReclamation(item) {
  const raw = item.statusUi || item.status || 'submitted';
  const statusKey =
    raw === 'resolved'     ? 'accepted'  :
    raw === 'under_review' ? 'reviewed'  :
    raw === 'rejected'     ? 'rejected'  : 'pending';

  const adminResponse = item.adminResponse?.message
    ? { ...item.adminResponse,
        message: cleanReplyText(item.adminResponse.message) }
    : item.reponse_admin
      ? { message:  cleanReplyText(item.reponse_admin),
          date:     item.reponse_admin_date }
      : null;

  return { ...item, statusKey, adminResponse };
}
\end{lstlisting}

\subsection{Formulaire de Soumission}

\begin{lstlisting}[language=JavaScript,
  caption=RequestsPage.jsx — Soumission d'une réclamation]
const handleSubmitReclamation = async (e) => {
  e.preventDefault();
  setSubmitting(true);
  try {
    await api.post('/requests/reclamations', {
      sujet:   form.sujet,
      message: form.message,
      type:    'reclamation',
    });
    setSuccess('Votre reclamation a ete soumise avec succes.');
    loadMyRequests();
    resetForm();
  } catch (err) {
    setError(
      err.response?.data?.error?.message || 'Erreur lors de la soumission.'
    );
  } finally {
    setSubmitting(false);
  }
};
\end{lstlisting}

\subsection{Décision Administrative (Frontend)}

\begin{lstlisting}[language=JavaScript,
  caption=AdminRequestsPage.jsx — Envoi d'une décision]
const submitDecision = async () => {
  try {
    const { data } = await api.post(
      `/requests/admin/${selectedReq.id}/decide`,
      { action: decision, reponse: replyText || undefined }
    );
    // Mise a jour optimiste de la liste locale
    setRequests(prev => prev.map(r =>
      r.id === selectedReq.id
        ? { ...r, status: data.status,
                  adminResponse: data.adminResponse,
                  updatedAt: data.updatedAt }
        : r
    ));
    closeModal();
  } catch (err) {
    const msg = err.response?.data?.error?.message || '';
    if (msg.includes('already') || msg.includes('deja'))
      alert('Cette demande a deja ete traitee. Rafraichissez la page.');
  }
};
\end{lstlisting}

\subsection{Composant de Conversation (Fil de Dialogue)}

\begin{lstlisting}[language=JavaScript,
  caption=RequestConversation.jsx — Badge statut et réponse admin]
const STATUS_META = {
  pending:  { label: 'En attente', className: 'bg-yellow-500/20 text-yellow-300' },
  reviewed: { label: 'En cours',   className: 'bg-blue-500/20  text-blue-300'   },
  accepted: { label: 'Acceptee',   className: 'bg-green-500/20 text-green-300'  },
  rejected: { label: 'Rejetee',    className: 'bg-red-500/20   text-red-300'    },
};

// Badge de statut
<span className={`px-2 py-1 rounded-full text-xs font-semibold
  ${STATUS_META[statusKey].className}`}>
  {STATUS_META[statusKey].label}
</span>

// Reponse admin ou etat vide
{adminResponse?.message ? (
  <div className="bg-blue-500/10 border border-blue-500/20 rounded-lg p-4">
    <p className="text-sm text-gray-300">{adminResponse.message}</p>
    {adminResponse.date && (
      <p className="text-xs text-gray-500 mt-1">
        {new Date(adminResponse.date).toLocaleDateString('fr-FR')}
      </p>
    )}
  </div>
) : (
  <p className="text-center text-gray-500 text-sm py-4">
    Aucune reponse de l'administration pour le moment.
  </p>
)}
\end{lstlisting}

\subsection{Formulaire Visiteur (Page Publique)}

\begin{lstlisting}[language=JavaScript,
  caption=PublicReclamationForm.jsx — Soumission sans authentification]
const handleGuestSubmit = async (e) => {
  e.preventDefault();
  try {
    await api.post('/requests/guest/reclamations', {
      prenom:  form.prenom,
      nom:     form.nom,
      email:   form.email,
      sujet:   form.sujet,
      message: form.message,
    });
    setConfirmed(true); // affiche le message de confirmation
  } catch {
    setError('Une erreur est survenue. Veuillez reessayer.');
  }
};
\end{lstlisting}

% ────────────────────────────────────────────────────────────
\section{Interfaces Utilisateur}
% ────────────────────────────────────────────────────────────

\subsection{Tableau de Bord Étudiant}

\begin{figure}[H]
\centering
\begin{tikzpicture}
\draw[IBNBlue,thick,rounded corners=6pt,fill=LightBlue!30]
  (0,0) rectangle (13,7);
% En-tete
\draw[IBNBlue,fill=IBNBlue,rounded corners=6pt 6pt 0 0]
  (0,6.3) rectangle (13,7);
\node[font=\small\bfseries\color{white}] at (6.5,6.65)
  {Mes Reclamations \& Justifications};
% Cartes statistiques
\foreach \x/\n/\lbl/\col in {
  2/3/{En attente}/yellow!30,
  5/1/{En cours}/blue!20,
  8/2/{Traitees}/green!20,
  11/6/Total/gray!20}{
  \draw[\col,draw=gray!40,rounded corners=4pt,fill=\col]
    (\x-1.3,5.0) rectangle (\x+1.3,6.0);
  \node[font=\large\bfseries\color=IBNBlue] at (\x,5.65) {\n};
  \node[font=\scriptsize\color{AccentGray}] at (\x,5.15) {\lbl};
}
% En-tete tableau
\draw[gray!30,fill=gray!10] (0.3,3.9) rectangle (12.7,4.6);
\node[font=\scriptsize\bfseries] at (2.5,4.25)  {Type};
\node[font=\scriptsize\bfseries] at (5.5,4.25)  {Sujet};
\node[font=\scriptsize\bfseries] at (8.5,4.25)  {Date};
\node[font=\scriptsize\bfseries] at (11.0,4.25) {Statut};
% Lignes
\foreach \y/\type/\sujet/\date/\stat/\scol in {
  3.4/Reclamation/{Note incorrecte...}/{12 avr.}/{En attente}/yellow!40,
  2.7/Justification/{Absence medicale}/{10 avr.}/Acceptee/green!30,
  2.0/Reclamation/{Erreur emploi du temps}/{5 avr.}/Rejetee/red!25}{
  \draw[gray!20] (0.3,\y-0.3) -- (12.7,\y-0.3);
  \node[font=\scriptsize] at (2.5,\y)  {\type};
  \node[font=\scriptsize] at (5.5,\y)  {\sujet};
  \node[font=\scriptsize] at (8.5,\y)  {\date};
  \draw[\scol,rounded corners=8pt,fill=\scol]
    (9.5,\y-0.18) rectangle (12.5,\y+0.18);
  \node[font=\tiny\bfseries] at (11.0,\y) {\stat};
}
\end{tikzpicture}
\caption{Interface étudiant — Liste des demandes avec statuts colorés}
\label{fig:student_ui}
\end{figure}

\subsection{Boîte de Réception Administrateur}

\begin{figure}[H]
\centering
\begin{tikzpicture}
\draw[IBNBlue,thick,rounded corners=6pt,fill=LightBlue!20]
  (0,0) rectangle (13,8);
% En-tete
\draw[IBNBlue,fill=IBNBlue,rounded corners=6pt 6pt 0 0]
  (0,7.2) rectangle (13,8);
\node[font=\small\bfseries\color{white}] at (6.5,7.6)
  {Gestion des Reclamations --- Administration};
% Onglets
\foreach \x/\lbl/\sel in {
  1.8/{Toutes}/1, 4.5/{Reclamations}/0,
  7.5/{Justifications}/0, 10.8/{Visiteurs}/0}{
  \draw[IBNBlue!\sel 00!gray!40,
        fill=IBNBlue!\sel 15!white,rounded corners=4pt]
    (\x-1.3,6.5) rectangle (\x+1.3,7.0);
  \node[font=\scriptsize\bfseries] at (\x,6.75) {\lbl};
}
% Stats
\foreach \x/\n/\lbl in {
  2/12/Total, 5/5/Nouvelles, 8/4/{En cours}, 11/3/Traitees}{
  \draw[IBNBlue!10,draw=IBNBlue!40,rounded corners=4pt]
    (\x-1.3,5.2) rectangle (\x+1.3,6.2);
  \node[font=\large\bfseries\color{IBNBlue}] at (\x,5.8) {\n};
  \node[font=\scriptsize\color{AccentGray}]  at (\x,5.35){\lbl};
}
% Liste
\draw[gray!15,fill=white,rounded corners=4pt]
  (0.3,0.3) rectangle (12.7,4.9);
\draw[gray!30,fill=gray!10] (0.3,4.3) rectangle (12.7,4.9);
\node[font=\scriptsize\bfseries] at (2.5,4.6)  {Etudiant};
\node[font=\scriptsize\bfseries] at (5.5,4.6)  {Type};
\node[font=\scriptsize\bfseries] at (8.5,4.6)  {Statut};
\node[font=\scriptsize\bfseries] at (11.0,4.6) {Action};
\foreach \y/\nom/\type/\stat/\scol/\btn in {
  3.8/{Fatima B.}/Reclamation/{En attente}/yellow!30/Traiter,
  3.1/{Karim M.}/Justification/{En cours}/blue!20/Voir,
  2.4/{Sara H.}/Reclamation/Acceptee/green!20/{---},
  1.7/{Ahmed L.}/Justification/Rejetee/red!20/{---},
  1.0/Visiteur/Reclamation/{En attente}/yellow!30/Traiter}{
  \draw[gray!15] (0.3,\y-0.3) -- (12.7,\y-0.3);
  \node[font=\scriptsize] at (2.5,\y) {\nom};
  \node[font=\scriptsize] at (5.5,\y) {\type};
  \draw[\scol,rounded corners=6pt,fill=\scol]
    (7.2,\y-0.18) rectangle (9.8,\y+0.18);
  \node[font=\tiny\bfseries] at (8.5,\y) {\stat};
  \draw[IBNBlue!20,fill=IBNBlue!15,rounded corners=4pt]
    (9.9,\y-0.18) rectangle (12.1,\y+0.18);
  \node[font=\tiny\bfseries\color{IBNBlue}] at (11.0,\y) {\btn};
}
\end{tikzpicture}
\caption{Boîte de réception administrateur avec filtres et actions}
\label{fig:admin_ui}
\end{figure}

% ────────────────────────────────────────────────────────────
\section{Tests et Validation}
% ────────────────────────────────────────────────────────────

\subsection{Scénarios de Tests Fonctionnels}

Les scénarios suivants ont été exécutés pour valider l'ensemble des
fonctionnalités du module~:

\begin{table}[H]
\centering
\caption{Résultats des tests fonctionnels — Réclamations \& Justifications}
\begin{tabularx}{\linewidth}{@{}lLll@{}}
\toprule
\textbf{ID} & \textbf{Scénario} & \textbf{Résultat attendu} & \textbf{Statut} \\
\midrule
TF-01 & Étudiant soumet une réclamation &
  Créée, statut \texttt{submitted} & PASS \\
TF-02 & Étudiant soumet une justification avec date d'absence &
  Créée avec \texttt{dateAbsence} & PASS \\
TF-03 & Soumission avec message vide &
  Erreur \texttt{400 MISSING\_FIELDS} & PASS \\
TF-04 & Champs nom/email pré-remplis depuis le compte &
  Non modifiables par l'étudiant & PASS \\
TF-05 & Étudiant consulte ses demandes &
  Liste filtrée à ses propres requêtes & PASS \\
TF-06 & Étudiant tente de voir la demande d'un autre &
  Accès refusé \texttt{403} & PASS \\
TF-07 & Admin consulte la boîte de réception &
  Toutes les demandes affichées & PASS \\
TF-08 & Admin approuve une réclamation (avec réponse) &
  Statut \texttt{resolved}, étudiant notifié & PASS \\
TF-09 & Admin approuve sans rédiger de réponse &
  Accepté, étudiant notifié & PASS \\
TF-10 & Admin rejette une justification &
  Statut \texttt{rejected}, étudiant notifié & PASS \\
TF-11 & Admin retraite une demande déjà traitée &
  Erreur \texttt{409 ALREADY\_REVIEWED} & PASS \\
TF-12 & Visiteur soumet une réclamation (sans compte) &
  Sauvegardée dans \texttt{GuestReclamation} & PASS \\
TF-13 & Visiteur tente une justification &
  Option indisponible (page publique) & PASS \\
TF-14 & Étudiant reçoit notification après décision &
  Notification visible dans le dashboard & PASS \\
\bottomrule
\end{tabularx}
\end{table}

\subsection{Tests de Sécurité}

\begin{table}[H]
\centering
\caption{Tests de sécurité — Module Réclamations \& Justifications}
\begin{tabularx}{\linewidth}{@{}lLl@{}}
\toprule
\textbf{Test} & \textbf{Description} & \textbf{Résultat} \\
\midrule
Accès sans token &
  Requête sur \texttt{/my-reclamations} sans JWT &
  401 TOKEN\_MISSING \\
Rôle incorrect &
  Étudiant accède à \texttt{/admin/inbox} &
  403 FORBIDDEN \\
Injection via message &
  Payload SQL/XSS dans le champ message &
  Neutralisé (Prisma) \\
Double décision &
  Admin retraite une demande déjà rejetée &
  409 ALREADY\_REVIEWED \\
\bottomrule
\end{tabularx}
\end{table}

\subsection{Taux de Couverture}

\begin{figure}[H]
\centering
\begin{tikzpicture}
\begin{axis}[
  xbar,
  width=10cm, height=5.5cm,
  bar width=0.5cm,
  xlabel={Couverture (\%)},
  symbolic y coords={Notification,Visiteur,Admin,Etudiant,Routes},
  ytick=data,
  xmin=0, xmax=100,
  xmajorgrids=true,
  grid style={dashed,gray!30},
  nodes near coords,
  nodes near coords align={horizontal},
  every node near coord/.append style={font=\small},
]
\addplot[fill=IBNBlue!60,draw=IBNBlue] coordinates {
  (88,Notification)
  (92,Visiteur)
  (95,Admin)
  (94,Etudiant)
  (100,Routes)
};
\end{axis}
\end{tikzpicture}
\caption{Couverture des tests par fonctionnalité}
\label{fig:rj_coverage}
\end{figure}

% ────────────────────────────────────────────────────────────
\section{Conclusion}
% ────────────────────────────────────────────────────────────

Ce chapitre a présenté l'implémentation complète du module
Réclamations \& Justifications~: modèle de données Prisma, routes API
sécurisées par JWT et RBAC, contrôleur avec gestion du cycle de vie des
statuts, service de notification automatique, et interfaces React adaptées
aux trois profils utilisateurs (étudiant, administrateur, visiteur).

La campagne de tests fonctionnels et de sécurité confirme que toutes les
fonctionnalités sont opérationnelles~: les étudiants soumettent et suivent
leurs demandes, les administrateurs disposent d'une boîte de réception
centralisée avec décisions traçables, et les visiteurs peuvent signaler des
problèmes sans créer de compte. Ce module élimine entièrement le processus
papier antérieur et améliore significativement la communication entre les
étudiants et l'administration universitaire.
